% \subsection{Black Speculative Traditions and Sonic Fiction}
\subsection{Afrofuturism as Aesthetic and Methodology}
Afrofuturism has emerged as a critical site for examining the intersections of race, culture, and technology. Ytasha Womack's reading of the term is Afrofuturism is as "an intersection of imagination, technology, the future, and liberation"\parencite[p. 9]{Womack_2013}. This widely cited formulation highlights Afrofuturism’s dual role as both artistic aesthetic and methodological orientation. For Womack, Afrofuturism encompasses a broad range of practices. These are anywhere from the literary, to the musical, the visual, and the performative. All these many "disciplines" within Afrofuturism deploy speculative imagination to reconfigure how Black culture relates to technology and futurity. She surveys a genealogy that includes musicians like Sun Ra, novelists like Octavia Butler, and contemporary artists like Janelle Monáe, arguing that each exemplifies how speculative creativity has been a persistent mode of Black survival and liberation. Afrofuturism is a way of thinking and knowing that treats cultural memory and futurist speculation as mutually reinforcing forces.  

Alondra Nelson, in her introduction to the 2002 \emph{Social Text} special issue, offers a complementary but distinct framing. She characterizes Afrofuturism as a means to "excavate and create original narratives of identity, technology, and the future and offer critiques of the promises of prevailing theories of technoculture." \parencite[p. 9]{Nelson_2002}. Nelson stresses that Afrofuturism should not be reduced to a canon of texts or artworks; instead, it should be understood as an evolving discourse that both documents and critiques the entanglements of race and technology. In her insistence on viewing Afrofuturism as an open conversation over that of a closed tradition, she positions it as at once a methodological resource and an aesthetic practice..  

Reading Nelson and Womack together raises important questions. Is Afrofuturism primarily an aesthetic framework that emerges from cultural production, or is it a critical methodology for interrogating technology? Womack leans toward emphasizing cultural archive and creative expression, while Nelson highlights scholarly and methodological potential. Yet rather than representing a contradiction, this dual emphasis reveals Afrofuturism’s true strength: it is simultaneously descriptive and prescriptive, drawing from cultural practice while also offering critical tools for future-oriented analysis.  

Another area of convergence is their treatment of temporality. Womack emphasizes that Afrofuturism draws on African diasporic legacies and memories as a means of inspiring speculative visions of the future \parencite{Womack_2013}. Nelson similarly observes, through investigating the work of Ishmael Reed, that Afrofuturist practices (and other futurisms) often spotlight histories and traditions of the displaced by re-situating them in the presence of current technological imaginaries \parencite[8]{Nelson_2002}. Both suggest that Afrofuturism collapses the binary of past and future, insisting instead on their entanglement in the present. This raises methodological implications: how might technological design processes integrate memory and futurity simultaneously? What does it mean to construct non-linear futures that refuse dominant narratives of progress?  

For this project, Afrofuturism, as articulated by Nelson and Womack and expanded by critical scholarship, offers four key contributions. First, it demonstrates that imagination itself can be treated as method, making speculation a legitimate mode of knowledge production. Second, it positions Black cultural traditions as central to rethinking technology. Third, it foregrounds technology as a contested site, one that can either reproduce exclusion or be re-engineered for liberation. Fourth, it generates liberation technology design principles that translate critique into actionable frameworks for Black autonomy. In sum, Afrofuturism supplies both the cultural archive and the methodological framework necessary for grounding Sonic Speculocultural Technopoiesis in practices of Black speculative thought.

\subsection{Sonic Fiction and Speculative Perception}
Kodwo Eshun extends Afrofuturist discourse into the explicitly sonic domain. In \emph{More Brilliant Than the Sun: Adventures in Sonic Fiction}, he introduces the concept of its subtitle, \textbf{sonic fiction}, proposing that Black electronic music—funk, dub, hip hop, techno—should be read as speculative technology the moves far beyond text, in its own right.

\begin{quotation}
	"The reason I don't talk about the literary is that there's just no need to, what with thinking about amplification and the sensory environment of amplification, of loudness in itself, the sensory impact of volume, the sensory impact of repetition, of broadcasting, all these things." \parencite [p. 183]{Eshun_1999}
\end{quotation}

 Unlike written fictions, which projects futures through literary narrative, sonic fictions enact futures directly through sound. Eshun insists that music/sound itself is a speculative apparatus meant for the active reconfiguration of identity, temporality, and perception. One meant to eclipse reality and "put out the sun. \parencite[p. 4]{Eshun_1999}"   

Black speculative thought has long been mediated sonically, through modes of listening, performance, and noise that trouble the boundaries between speech, song, and world-making.  As Alexander G. Weheliye argues, “the sonic remains an important zone from and through which to theorize the fundamentality of Afro-diasporic formations to the currents of Western modernity, since this field remains the principal modality in which Afro-diasporic cultures have been articulated...” \parencite[p.5]{Weheliye_2005}. Sound, in this sense, is fundamentally ontological: a means through which Black life becomes audible and knowable to itself against the muting of colonial modernity and its many technological fixes and interventions. The groove, the remix, and the recorded voice each constitute what Weheliye terms \emph{phonographic} inscription, sites where memory, identity, and futurity are encoded through vibration, as opposed text (or dominantly visual media). These inscriptions enact a counter-archive to Eurocentric epistemologies by privileging the performative and the improvisatory as modes of knowledge. 

Recent scholarship extends this understanding through the language of Black noise and sonic invention. Dorottya Mozes identifies \emph{black noise} as the sonic modality of "protest against violent domination,” a field where listening becomes a practice of survival that “fosters alternative forms of relationality and spatiality” \parencite[167-172]{Mozes_2022}.  Such sonic experiments rehearse the conditions of Black life, transforming the so-called excesses of noise or experience-junk into expressive acts of world-making.  Julia Reade likewise describes these performances as “breaks” where "improvised ruptures in Black performances, expose a liminal space brimming with generative and ontological potential..."  \parencite[3]{Reade_2023}.  From this perspective, the aesthetics of Afrofuturism and sonic fiction become practices for the conversion of the sonic residue of racialized constraint into the improvisatory space. In this place, listening emerges as a critical and creative method as a part of sonic thinking, where the future is composed in vibration father than text.       

Eshun provides vivid examples of this speculative function. When describing Cybotron’s "Techno City," he calls it "...electronic fiction, with frequencies fictionalized, synthesized, and organized into escape routes." \parencite[p. 103]{Eshun_1999}. Such sonic fictions are heard as real interventions that carve perceptual exits from the “Cartesian prison” of your senses on pure and proper values of conventional musicality. In these instances, sound functions as a mode of engineering: it reorganizes how listeners inhabit time, space, and social reality. Eshun describes this as seizing "the means of perception" \parencite[p. 32]{Eshun_1999}, underscoring that sonic practices alter both aurality and perception.  

One striking aspect of Eshun’s analysis is his attention to temporality. In this, we see the influence of Afrofuturist thought on Eshun's work. Without explicitly naming it as such, he shows how repetition, delay, looping, and acceleration in several nascent practices, makes of Black music operate as strategies for deftly manipulating temporal experience. Dub’s recursive echoes or Detroit Techno’s temporal elasticity flourish via speculative manipulations of time that enable listeners to inhabit the non-linear. It casts us deep into "chronography," the age in which we inhabit time. \parencite[p. 101]{Eshun_1999} This implicit theorization of time, too, resonates with Womack's claim that Afrofuturism, as a critical theory, is as much about varied intertwinings with imagination, as it is about future projection and cultural critiques. \parencite[p. 9]{Womack_2013}. Both accounts suggest that Black speculative practices produce futures that are recursive, layered, and entangled with the baggage of histories.  

The technological dimension of Eshun’s argument is equally important. Instruments such as turntables, samplers, and synthesizers are recast as speculative machines. As apparatuses that fictionalize sound into future-space. This often stands in stark contrast to media hierarchies that privilege literature or visual culture as the primary sites of futurist imagination. By placing sonic practice at the center of speculative world-making, Eshun reframes technology itself; in a way that punches past the all to obvious hardware or software qualities, to present embodied cultural logics encoded through sound.  

As it relates to \emph{Sonic Speculocultural Technopoiesis}, Eshun's concept of sonic fiction presents Black sonic principles as organizational factors for perception and interaction. If we hope to treat sonic fiction as speculative technology we must begin by imagining systems that embody its dialectics.  
  


